% ============================================================================
% CV Youssef EL KATTOUFI - candidature Bell 2027, novembre 2026
% LuaLaTeX recommande (Overleaf > Menu > Compilateur > LuaLaTeX)
% Mise a jour personnalisee a partir du CV source et de l'offre Bell 432858.
% Les objectifs sont nommes comme tels ; aucune performance non mesuree n'est presentee comme acquise.
% Identite : bleu / rouge inspires du logo Sorbonne Universite ; gris du
%           langage graphique Polytech. Aucune reproduction de logo officiel.
% Version sans illustrations dans les marges, avec identite bleu / rouge.
% ============================================================================
\documentclass[9pt,a4paper]{extarticle}
\usepackage[a4paper,top=0.95cm,bottom=0.82cm,left=1.24cm,right=1.24cm]{geometry}
\usepackage{fontspec}
\setmainfont{Carlito}[BoldFont={Carlito Bold}, ItalicFont={Carlito Italic},Ligatures=NoCommon]
\usepackage[french]{babel}
\usepackage{microtype}
\usepackage{xcolor}
\usepackage{tabularx}
\usepackage{array}
\usepackage{enumitem}
\usepackage{titlesec}
\usepackage[hidelinks]{hyperref}
\hypersetup{pdftitle={CV - Youssef EL KATTOUFI - Electronique et Informatique},pdfauthor={Youssef EL KATTOUFI},pdfsubject={Candidature au programme de stages 2027 en innovation technologique chez Bell},colorlinks=true,urlcolor=bleu}
\definecolor{bleu}{HTML}{263068}
\definecolor{rouge}{HTML}{E6332A}
\definecolor{gris}{HTML}{9EA6B6}
\definecolor{anthracite}{HTML}{222936}
\definecolor{secondaire}{HTML}{475269}
\definecolor{fond}{HTML}{F0F2F7}
\color{anthracite}
\pagestyle{empty}
\setlength{\parindent}{0pt}
\setlength{\parskip}{0pt}
\setlength{\emergencystretch}{1em}
\renewcommand{\arraystretch}{0.98}
\newcommand{\stage}{STAGE BELL 2027 -- DISPONIBLE DÈS AVRIL 2027}

% --- Composants typographiques ------------------------------------------------
\titleformat{\section}{\fontsize{9.85}{11}\bfseries\color{bleu}}{}{0pt}{\MakeUppercase}[{\vspace{-2pt}\color{rouge}\titlerule[.66pt]}]
\titlespacing*{\section}{0pt}{7.0pt}{3.1pt}
\setlist[itemize]{leftmargin=11pt,topsep=.9pt,itemsep=.55pt,parsep=0pt,partopsep=0pt,label={\color{rouge}\raisebox{.4ex}{\tiny$\bullet$}}}
\newcommand{\entry}[4]{%
\begin{tabularx}{\linewidth}{@{}>{\raggedright\arraybackslash}X>{\raggedleft\arraybackslash}p{3.08cm}@{}}
 {\bfseries\color{bleu}#1} & {\fontsize{8.5}{9.3}\selectfont\color{secondaire}#2} \\[-2pt]
 {\fontsize{8.5}{9.6}\selectfont\color{secondaire}#3} & {\fontsize{8.5}{9.3}\selectfont\color{secondaire}#4}
\end{tabularx}\vspace{-3.4pt}\par}
\newcommand{\gapentry}{\vspace{3.0pt}}
\newcommand{\comp}[2]{\textbf{\color{bleu}#1} \hspace{2pt} #2\par}


\begin{document}
% Profil et coordonnees ; couleurs et structure d'origine conservees
\begin{minipage}[t]{.635\linewidth}
  {\fontsize{21.2}{22.3}\selectfont\bfseries\color{bleu}Youssef EL KATTOUFI}\par\vspace{1.5pt}
  {\fontsize{10.2}{11.5}\selectfont\bfseries\color{rouge}ÉLÈVE-INGÉNIEUR -- ÉLECTRONIQUE ET INFORMATIQUE}\par\vspace{1.1pt}
  {\fontsize{9.15}{10}\selectfont\bfseries\color{bleu}\stage}\par
\end{minipage}\hfill
\begin{minipage}[t]{.355\linewidth}\raggedleft
  {\fontsize{8.7}{10.2}\selectfont
  Choisy-le-Roi, France\\
  \href{tel:+33651654237}{+33 6 51 65 42 37}\\
  \href{mailto:elkattoufi.youssef@hotmail.com}{elkattoufi.youssef@hotmail.com}\\
  \href{https://www.linkedin.com/in/youssef-el-kattoufi-5a04a0282/}{LinkedIn} \quad\textcolor{gris}{|}\quad \href{https://github.com/Darkzveller}{GitHub} }
\end{minipage}
\vspace{4.1pt}
\textcolor{rouge}{\rule{1.35cm}{1.75pt}}\hspace{1pt}\textcolor{gris!65}{\rule{\dimexpr\linewidth-1.45cm\relax}{.6pt}}\par\vspace{3.0pt}
{\fontsize{9.0}{10.7}\selectfont
Élève-ingénieur en Électronique et Informatique à \textbf{Polytech Sorbonne}, je recherche un \textbf{stage dès avril 2027}. Mes expériences m'ont appris à \textbf{développer et tester} des systèmes, à analyser des problèmes concrets et à \textbf{travailler avec plusieurs équipes}. Je souhaite mettre ces acquis en pratique dans les domaines du logiciel, des réseaux ou des systèmes connectés.\par}

\section{Expériences professionnelles}
\entry{Surveillant sauveteur aquatique}{Juil. -- août 2026}{Centre aquatique de Villecresnes -- GPSEA}{Villecresnes}
\begin{itemize}
  \item \textbf{Contexte :} assurer la sécurité d'un établissement pouvant accueillir jusqu'à \textbf{1\,000 personnes}.
  \item \textbf{Missions :} surveillance et prévention des risques en coordination avec les autres sauveteurs, communication avec le public et interventions de premiers secours.
\end{itemize}
\gapentry
\entry{Apprenti ingénieur -- Génie électrique}{Oct. 2024 -- août 2025}{Syctom -- sites Isséane et Saint-Ouen}{Île-de-France}
\begin{itemize}
  \item \textbf{Sujet :} modernisation d'installations techniques en environnement industriel, avec des échanges entre exploitants, prestataires et équipes de maintenance.
  \item \textbf{Missions et contributions :} analyse des besoins et rédaction de pièces techniques pour le renouvellement de la \textbf{GTB} ; suivi de la modernisation des liaisons de sûreté en \textbf{fibre optique}, commandes et coordination de chantier. Étude de coût et d'amortissement de \textbf{6 bornes de recharge}.
\end{itemize}
\gapentry
\entry{Stagiaire technicien électronicien supérieur}{Avr. -- juin 2024}{Cobham Aerospace Communications}{France}
\begin{itemize}
  \item \textbf{Sujet :} fiabiliser et accélérer le contrôle de cartes électroniques.
  \item \textbf{Missions :} développement en \textbf{C++} de séquences de tests automatisés, création d'une interface de commande, réalisation d'essais et diagnostic des anomalies.
  \item \textbf{Livrables et outils :} banc de test, documentation technique et simulations sous \textbf{LTspice}. \textbf{Objectif :} estimation de la réduction de la durée de test de \textbf{30\,\%}.
\end{itemize}

\section{Projets techniques \& associatifs}
\entry{Projet IoT -- Acquisition et partage de mesures}{2026}{Polytech Sorbonne -- ESP32 et capteur de température/humidité}{Universitaire}
\begin{itemize}
  \item \textbf{Sujet :} rendre accessibles à distance les mesures d'un capteur connecté, pour pouvoir les consulter et les exploiter.
  \item \textbf{Missions et résultat :} acquisition des données sur \textbf{ESP32}, transmission via \textbf{MQTT} et \textbf{HTTP REST}, puis création de flux \textbf{Node-RED} pour le suivi des mesures.
\end{itemize}
\gapentry
\entry{Traitement du signal -- Réception et décodage DCF77}{2025 -- 2026}{Polytech Sorbonne -- radio à 77,5 kHz}{Universitaire}
\begin{itemize}
  \item \textbf{Sujet :} retrouver les informations binaires d'un signal radio \textbf{DCF77} bruité à partir de mesures à l'oscilloscope et d'enregistrements WebSDR.
  \item \textbf{Missions :} traitement sous \textbf{MATLAB} par filtre RII d'ordre 8, \textbf{filtfilt}, extraction de l'enveloppe par filtre de Butterworth et seuillage adaptatif.
  \item \textbf{Résultat :} reconstruction du signal binaire et identification des bits selon des impulsions de \textbf{100 et 200 ms}.
\end{itemize}
\gapentry
\entry{Coupe de France de robotique -- Responsable électronique / asservissement}{2024 -- 2025}{Club CRAC -- robot autonome développé en équipe}{Associatif}
\begin{itemize}
  \item \textbf{Sujet :} intégrer la motorisation, les cartes électroniques et le logiciel d'un robot destiné à la Coupe de France de robotique.
  \item \textbf{Missions :} développement du socle \textbf{C++} sur \textbf{ESP32/STM32}, échanges \textbf{CAN}, pilotage des actionneurs, interface de tests et mises à jour \textbf{OTA}.
  \item \textbf{Réalisations :} commande de \textbf{2 moteurs}, odométrie et asservissements \textbf{PID} en vitesse/position ; intégration avec l'équipe. \textbf{Outils :} FreeRTOS, Git, PlatformIO.
\end{itemize}
\gapentry
\entry{Trottinette à double motorisation -- moteur sans terres rares}{2024 -- 2025}{IUT de Cachan -- électronique de puissance et conception}{Universitaire}
\begin{itemize}
  \item \textbf{Sujet :} concevoir le pilotage d'un \textbf{moteur AC sans terres rares} associé à un moteur \textbf{DC} d'assistance.
  \item \textbf{Missions :} réalisation d'une \textbf{carte PCB} et développement de commandes \textbf{STM32/ESP32} ; suivi de la température, des phases, de la vitesse et de la consommation électrique.
  \item \textbf{Réalisations :} rebobinage moteur, odométrie, conception \textbf{CAO 3D} et essais avec un \textbf{onduleur triphasé}.
\end{itemize}

\section{Formation}
\entry{Diplôme d'ingénieur -- Électronique et Informatique}{2025 -- 2028}{Polytech Sorbonne -- Sorbonne Université}{Paris}
\vspace{1pt}
\entry{BUT Génie électrique et informatique industrielle -- option Électricité et maîtrise de l'énergie}{2022 -- 2025}{IUT de Cachan -- Université Paris-Saclay}{Cachan}

\section{Compétences}
{\fontsize{8.65}{10.05}\selectfont
\comp{Programmation :}{C, C++, Python, MATLAB.}
\comp{Systèmes et communication :}{STM32, ESP32, CAN, UART, SPI, I\textsuperscript{2}C, PID, FreeRTOS, MQTT, HTTP REST.}
\comp{Logiciels et conception :}{PlatformIO, Git, LTspice, Node-RED, KiCad, SolidWorks, STM32CubeIDE.}
}
\end{document}
